\section{Background}


\subsection{Rust for ML}

Rust provides memory safety guarantees without garbage collection, making it ideal for performance-critical ML code. The ownership model prevents data races in parallel computation, while zero-cost abstractions enable high-level APIs without runtime overhead.

Prior work on Rust ML frameworks includes tch-rs (PyTorch bindings)~\cite{tchrs2020} and burn~\cite{burn2023}. Our framework differs by prioritizing deployment efficiency and FHE compatibility over training performance.

\subsection{Fully Homomorphic Encryption}

FHE schemes like TFHE~\cite{tfhe2016} and CKKS~\cite{ckks2017} enable computation on encrypted data. However, FHE supports only polynomial operations efficiently; non-polynomial functions (ReLU, sigmoid, softmax) require approximation or bootstrapping.

\subsection{Quantization}

Model quantization reduces precision from FP32 to INT8/INT4, decreasing memory bandwidth and enabling vectorized computation. The GGML format~\cite{ggml2023} introduces k-quants (Q4\_K, Q5\_K, Q6\_K) that combine multiple quantization scales per block for improved accuracy.

